\documentclass[11pt]{article}
\usepackage[margin=0.9in]{geometry}
\usepackage{amsmath,amssymb,amsthm,mathrsfs}
\usepackage{graphicx}
\usepackage[hidelinks]{hyperref}
\usepackage{microtype}

\newtheorem*{theorem}{Theorem}
\newtheorem*{lemma}{Exact finite certificate lemma}
\newcommand{\E}{\mathscr E}
\newcommand{\C}{\mathscr C}
\newcommand{\one}{\mathbf 1}
\newcommand{\R}{\mathbb R}
\setlength{\emergencystretch}{3em}
\graphicspath{{figures/}}

\begin{document}

\begin{center}
{\Large ETF SOS Depth Is Not Determined by $(N,r)$}\\[0.45em]
{\large An exact counterexample to Question 4.1 of arXiv:1901.10697}\\[0.35em]
{\small Lane 19, agent\_lane\_19, model GPT5.6}
\end{center}

\paragraph{Status.}
\texttt{candidate\_full\_counterexample\_likely\_valid}.  The finite part of
the proof is certified with exact integer and rational arithmetic, rather than
floating-point SDP output.

\paragraph{Source question.}
Bandeira--Kunisky, arXiv:1901.10697, Question 4.1 on PDF page 12 asks for the
largest even $d$ such that the Gram matrix of a real ETF of $N$ vectors in
$\R^r$ belongs to the degree-$d$ generalized elliptope $\E_d^N$, and asks
whether this largest degree depends only on $N$ and $r$.

\begin{center}
\includegraphics[width=0.96\textwidth]{open_problem_crop.png}
\end{center}

\paragraph{Definitions.}
The cut polytope is
\[
 \C^N=\operatorname{conv}\{zz^\top:z\in\{\pm1\}^N\}.
\]
The generalized elliptopes form the decreasing SOS hierarchy
\[
 \E_2^N\supseteq\E_4^N\supseteq\cdots\supseteq
 \E_{N+\mathbf1\{N\text{ odd}\}}^N=\C^N.                       \tag{1}
\]
For a graph with adjacency matrix $A$, its Seidel matrix is
$S=\one\one^\top-I-2A$.

\paragraph{Proof intuition.}
Strongly regular graphs with parameters $(26,10,3,4)$ all have Seidel
spectrum $\{5^{(13)},-5^{(13)}\}$.  Consequently $I+S/5$ is always the Gram
matrix of a real $(26,13)$ ETF.  What varies with the switching class is the
population of sign vectors in the $+5$ eigenspace.  In the first example that
population is rich enough that its uniform covariance is the ETF Gram matrix.
In the second, there are only fourteen sign lines.  Positivity forces any cut
decomposition to use only those lines, but their unique affine representation
of the Gram matrix assigns one line negative mass.  This gives two identical
frame parameter pairs but different terminal SOS membership.

\begin{theorem}
There exist two real equiangular tight frames with the same parameters
$(N,r)=(26,13)$ whose Gram matrices have different largest generalized-
elliptope degrees.  One Gram matrix belongs to $\E_{26}^{26}=\C^{26}$, while
the other does not.  Hence the answer to the second sentence of Question 4.1
does not depend only on $N$ and $r$.
\end{theorem}

\paragraph{The two exact graphs.}
Let $G_c,G_n$ be the graphs encoded by the following graph6 strings,
respectively:
\begin{verbatim}
Y}qCCmMlDSIPRCXPKSQhHQooE^`E`pHQoQuHBRc_oo}EHSMILOkEWXB_
Y}qCCmMlDSIQRBWokcadAROKE]XEXPHecTJ_AlWGooveDEcISRDEX_p_
\end{verbatim}
They are the first two entries in Brendan McKay's catalogue of the ten
strongly regular graphs with parameters $(26,10,3,4)$.  Write $A_c,A_n$ for
their adjacency matrices and
\[
 S_\bullet=\one\one^\top-I-2A_\bullet,
 \qquad X_\bullet=I+\frac15S_\bullet.                            \tag{2}
\]

\begin{lemma}
The following identities hold exactly.
\begin{enumerate}
\item Both graphs are $\operatorname{SRG}(26,10,3,4)$ and
      $S_c^2=S_n^2=25I$.
\item Up to replacing $z$ by $-z$, the set
      $Z_c=\{z\in\{\pm1\}^{26}:S_cz=5z\}$ has 130 elements and
      \[
        5\sum_{z\in Z_c}zz^\top=130(5I+S_c).                    \tag{3}
      \]
\item The analogous set $Z_n$ has 14 elements.  The fourteen vectors
      \[
       F(z)=\bigl(1,(z_i z_j)_{0\le i<j<26}\bigr)\in\R^{326}
      \]
      are linearly independent.  The unique coefficients satisfying
      \[
       \bigl(1,((S_n)_{ij}/5)_{i<j}\bigr)
       =\sum_{z\in Z_n}w_zF(z)                                  \tag{4}
      \]
      are, in a deterministic lexicographic ordering,
      \[
        -\frac3{10},\underbrace{\frac1{10},\ldots,\frac1{10}}_{13\ \mathrm{times}}.
                                                                    \tag{5}
      \]
\end{enumerate}
\end{lemma}

\begin{proof}[Verification of the lemma]
The accompanying script \texttt{code/verify\_counterexample.py} decodes the
two embedded graph6 strings and checks the degree and common-neighbor
conditions, as well as $S^2=25I$, with integer arithmetic.

For each $S$, it computes a rational basis of $\ker(S-5I)$ and selects thirteen
pivot coordinates on which restriction is invertible.  Every sign vector in
the eigenspace is therefore uniquely determined by one of the $2^{13}$ sign
assignments on those coordinates.  A common-denominator integer
reconstruction exhausts these assignments, proving the counts 130 and 14.
It verifies (3) entry by entry.  In the second case it computes the exact
rational rank of the $14\times326$ feature matrix, solves a nonsingular
$14\times14$ subsystem, obtains (5), and verifies all 326 equations in (4).
The deterministic output is saved as
\texttt{code/exact\_certificate.json}.
\end{proof}

\begin{proof}[Proof of the theorem]
Because $S_\bullet^2=25I$ and $\operatorname{tr}(S_\bullet)=0$, the eigenvalues
of each $S_\bullet$ are $+5$ and $-5$, each with multiplicity 13.  Thus
\[
 P_{\bullet,+}=\frac12\left(I+\frac15S_\bullet\right)
\]
is the rank-13 orthogonal projector onto the $+5$ eigenspace and
$X_\bullet=2P_{\bullet,+}$.  In particular, $X_\bullet$ is positive
semidefinite, has diagonal one, has all off-diagonal entries of modulus $1/5$,
and has nonzero eigenvalue $2=26/13$.  It is therefore the Gram matrix of a
real ETF of 26 vectors in $\R^{13}$.

Dividing (3) by 650 gives
\[
 X_c=\frac1{130}\sum_{z\in Z_c}zz^\top,
\]
so $X_c\in\C^{26}=\E_{26}^{26}$.  Its largest hierarchy degree is 26.

Suppose instead that $X_n\in\C^{26}$.  Then
\[
 X_n=\sum_{z\in\{\pm1\}^{26}}\lambda_z zz^\top,
 \qquad \lambda_z\ge0,\quad\sum_z\lambda_z=1.                  \tag{6}
\]
If $u\in\ker X_n$ then
\[
 0=u^\top X_nu=\sum_z\lambda_z(u^\top z)^2.
\]
Hence every $z$ with $\lambda_z>0$ is orthogonal to $\ker X_n$, so it belongs
to $\operatorname{ran}X_n=\ker(S_n-5I)$.  Up to sign, (6) can therefore use
only the fourteen lines $Z_n$ in the lemma.  Taking diagonal and
off-diagonal entries in (6) produces exactly (4), so linear independence
forces the unique weights (5).  One is negative, contradicting (6).  Thus
$X_n\notin\C^{26}=\E_{26}^{26}$.

The hierarchy levels are even, so the largest degree for $X_n$ is at most 24,
whereas that for $X_c$ is 26.  Both frames nevertheless have the same
$(N,r)=(26,13)$, proving the claim.
\end{proof}

\paragraph{Scope and limitations.}
This completely answers the yes/no parameter-dependence sentence of Question
4.1 in the negative.  It does not determine the exact largest degree for
$X_n$; the source's degree-4 criterion gives $X_n\in\E_4^{26}$ because
$26<13\cdot14/2$, so its depth lies between 4 and 24.  Nor does the result
classify all real ETFs or all regular two-graph switching classes.

\paragraph{Novelty search.}
The run indexes had no exact or core-question match.  Bounded searches through
2026-08-13 used the exact Question 4.1 wording and the terms
``ETF generalized elliptope'', ``regular two-graph cut polytope'',
``SRG(26,10,3,4) cut polytope'', and ``(26,13) equiangular tight frame cut
polytope''.  They found the cited ETF/SOS papers, Laurent's hierarchy lower
bound, the graph catalogues, and general $(26,13)$ ETF constructions, but no
same-parameter counterexample or the exact affine certificate (5).

\paragraph{Human-review recommendation.}
Check first the kernel-support implication following (6), then independently
rerun the exact verifier and inspect its pivot-coordinate exhaustiveness and
the 326-equation affine certificate.  Subject to those checks, the
counterexample is ready for expert review.

\paragraph{References.}
\begin{itemize}
\item A. S. Bandeira and D. Kunisky, \emph{Sum-of-Squares Optimization and the
Sparsity Structure of Equiangular Tight Frames}, arXiv:1901.10697 (2019).
\item A. S. Bandeira and D. Kunisky, \emph{A Gramian Description of the
Degree 4 Generalized Elliptope}, arXiv:1812.11583 (2018).
\item M. Laurent, \emph{Lower Bound for the Number of Iterations in
Semidefinite Hierarchies for the Cut Polytope}, Math. Oper. Res. 28 (2003),
871--883.
\item B. D. McKay, \emph{Combinatorial Data: Strongly Regular Graphs},
\url{https://users.cecs.anu.edu.au/~bdm/data/graphs.html}, catalogue file
\texttt{sr261034.g6}, accessed 2026-08-13.
\end{itemize}

\end{document}

